\documentclass[11pt]{article}

\usepackage[margin=1in]{geometry}
\usepackage{amsmath,amssymb,bm}
\usepackage{booktabs}
\usepackage{hyperref}

\newcommand{\Rnav}{\mathcal{R}_{\mathrm{nav}}}
\newcommand{\sigmar}{\sigma_r}
\newcommand{\sigmav}{\sigma_v}
\newcommand{\GammaGeom}{\Gamma}
\newcommand{\xstate}{\bm{x}}
\newcommand{\zhat}{\hat{\bm{z}}}
\newcommand{\Phimat}{\bm{\Phi}}
\newcommand{\Qmat}{\bm{Q}}
\newcommand{\Rmat}{\bm{R}}
\newcommand{\Pmat}{\bm{P}}
\newcommand{\Kmat}{\bm{K}}
\newcommand{\Hmat}{\bm{H}}
\newcommand{\Imat}{\bm{I}}


\title{Mathematical Foundations of Risk-Aware CR3BP Navigation}
\author{Fabio Facin}
\date{\today}

\begin{document}

\maketitle

\begin{abstract}
This note explains the mathematical structure behind a risk-aware navigation
filter for cislunar spacecraft motion. The framework combines CR3BP dynamics,
a nonlinear range-bearing measurement model, an Extended Kalman Filter (EKF),
a scalar navigation-risk score, and adaptive noise rules. The CR3BP and EKF
provide the physics and estimation foundation, while the navigation-risk score
is the main research contribution.
\end{abstract}

\section{Physical Dynamics: CR3BP}

The planar spacecraft state is
\begin{equation}
  \xstate =
  \begin{bmatrix}
    x & y & \dot{x} & \dot{y}
  \end{bmatrix}^{T}.
\end{equation}
This state comes from classical mechanics in a rotating reference frame. The
model is the circular restricted three-body problem (CR3BP), where the Earth
and Moon are the massive bodies and the spacecraft has negligible mass.

The nondimensional rotating-frame equations are
\begin{align}
  \ddot{x}
  &=
  2\dot{y}
  + x
  - \frac{(1-\mu)(x+\mu)}{r_1^3}
  - \frac{\mu(x-1+\mu)}{r_2^3},
  \\
  \ddot{y}
  &=
  -2\dot{x}
  + y
  - \frac{(1-\mu)y}{r_1^3}
  - \frac{\mu y}{r_2^3},
\end{align}
with
\begin{align}
  r_1 &= \sqrt{(x+\mu)^2+y^2},\\
  r_2 &= \sqrt{(x-1+\mu)^2+y^2}.
\end{align}

The terms $2\dot{y}$ and $-2\dot{x}$ are Coriolis terms caused by the rotating
frame. The terms $x$ and $y$ come from the centrifugal potential,
\begin{equation}
  U_{\mathrm{eff}} = \frac{1}{2}(x^2+y^2).
\end{equation}
The gravitational terms come from Newtonian gravity in vector form:
\begin{equation}
  \bm{a} = -\frac{GM}{r^3}\bm{r}.
\end{equation}
The CR3BP is useful for risk-aware filtering because it is already nonlinear
and can be dynamically sensitive. That makes it a strong test case for EKF
performance and uncertainty growth.

\section{Measurement Model}

An observer located at $(x_o,y_o)$ measures range and bearing:
\begin{equation}
  \bm{z}
  =
  \begin{bmatrix}
    \rho\\
    \theta
  \end{bmatrix}
  =
  \begin{bmatrix}
    \sqrt{(x-x_o)^2+(y-y_o)^2}\\
    \tan^{-1}\!\left(\frac{y-y_o}{x-x_o}\right)
  \end{bmatrix}
  + \bm{\eta}.
\end{equation}
The range term is Euclidean distance, and the bearing term is the angle from
the observer to the spacecraft. Measurement noise is modeled as
\begin{equation}
  \bm{\eta} \sim \mathcal{N}(\bm{0},\Rmat),
\end{equation}
where $\Rmat$ represents sensor imperfections, photon noise, tracking error,
and environmental degradation.

This measurement model is nonlinear, so the filter must linearize it. That is
why the EKF is appropriate.

\section{Extended Kalman Filter}

\subsection{Prediction}

The EKF first predicts the spacecraft state through the nonlinear dynamics:
\begin{equation}
  \hat{\xstate}_{k|k-1}
  =
  f(\hat{\xstate}_{k-1|k-1}).
\end{equation}
In implementation, the function $f$ is obtained by numerically integrating the
CR3BP differential equations, for example with fourth-order Runge--Kutta.

The covariance prediction is
\begin{equation}
  \Pmat_{k|k-1}
  =
  \Phi_k \Pmat_{k-1|k-1}\Phi_k^{T}
  + \Qmat.
\end{equation}
Here $\Pmat$ is the uncertainty covariance, $\Qmat$ is process noise, and
\begin{equation}
  \Phi_k = \frac{\partial f}{\partial \xstate}
\end{equation}
is the dynamics Jacobian. The Jacobian comes from a first-order Taylor
expansion of the nonlinear dynamics.

\subsection{Update}

The Kalman gain is
\begin{equation}
  K_k
  =
  \Pmat H_k^{T}
  \left(
    H_k \Pmat H_k^{T} + \Rmat
  \right)^{-1},
\end{equation}
where $H_k$ is the measurement Jacobian. The gain comes from minimizing
mean-squared estimation error under Gaussian and locally linear assumptions.

The corrected estimate is
\begin{equation}
  \hat{\xstate}_{k|k}
  =
  \hat{\xstate}_{k|k-1}
  +
  K_k
  \left(
    \bm{z}_k - h(\hat{\xstate}_{k|k-1})
  \right).
\end{equation}

\begin{table}[h]
\centering
\begin{tabular}{ll}
\toprule
Term & Meaning\\
\midrule
$\Pmat$ & Current uncertainty\\
$\Rmat$ & Measurement noise\\
$\Qmat$ & Process or model noise\\
$K_k$ & How strongly the filter trusts the measurement\\
\bottomrule
\end{tabular}
\caption{Interpretation of core EKF terms.}
\end{table}

\section{Main Equation: Navigation Risk Score}

The central contribution is the navigation-risk score:
\begin{equation}
\boxed{
\Rnav(t)
=
w_r\frac{\sigmar(t)}{r_0}
+
w_v\frac{\sigmav(t)}{v_0}
+
w_g\frac{1}{\GammaGeom(t)+\epsilon}
+
w_mM(t)
+
w_lL(t)
}.
\end{equation}
This equation turns a high-dimensional navigation state into a single scalar
warning signal.

\subsection{Position Uncertainty}

The position uncertainty term is
\begin{equation}
  \sigmar(t) = \sqrt{\operatorname{tr}(\Pmat_{rr}(t))}.
\end{equation}
The block $\Pmat_{rr}$ is the position covariance submatrix. The trace sums
the total position variance, and the square root converts that variance scale
into a standard-deviation scale.

\subsection{Velocity Uncertainty}

The velocity uncertainty term is
\begin{equation}
  \sigmav(t) = \sqrt{\operatorname{tr}(\Pmat_{vv}(t))}.
\end{equation}
Velocity matters because orbital prediction is not controlled by position
alone. A velocity error can drive large future trajectory divergence.

\subsection{Geometry Term}

The geometry penalty is
\begin{equation}
  \frac{1}{\GammaGeom(t)+\epsilon}.
\end{equation}
The quantity $\GammaGeom(t)$ measures measurement geometry quality. Good
geometry gives a large $\GammaGeom(t)$ and therefore low risk. Bad geometry
gives a small $\GammaGeom(t)$ and therefore high risk. The small positive
constant $\epsilon$ prevents division by zero.

This is similar in spirit to GPS dilution of precision and to observability
metrics from estimation theory.

\subsection{Missed Measurements}

The missed-measurement term is
\begin{equation}
  M(t)
  =
  \begin{cases}
    1, & \text{measurement missed},\\
    0, & \text{measurement received}.
  \end{cases}
\end{equation}
When a measurement is missed, the filter performs prediction without a
corrective observation. In nonlinear orbital dynamics, that usually means
uncertainty grows.

\subsection{Lighting Penalty}

The lighting penalty $L(t)$ adds aerospace realism. Optical tracking can
degrade during poor illumination, eclipse, glare, low contrast, or other
visibility problems.

\subsection{Weights}

The weights $w_r,w_v,w_g,w_m,w_l$ are tuning parameters, not physical
constants. They define how strongly each risk source contributes to the final
score. Conceptually, the score behaves like a cost function or a Lyapunov-like
summary of the filter's uncertainty condition.

\section{Adaptive Rules}

The navigation-risk score can modify the EKF noise assumptions:
\begin{align}
  \Rmat(t)
  &=
  \Rmat_0
  \left[
    1+\beta\,\operatorname{clip}(\Rnav(t))
  \right],
  \\
  \Qmat(t)
  &=
  \Qmat_0
  \left[
    1+\alpha\,\operatorname{clip}(\Rnav(t))
  \right].
\end{align}
These are adaptive filtering rules, not standard EKF equations. When risk
increases, the filter inflates its measurement noise or process noise model.
That makes the estimator more cautious under stressed navigation conditions.

\section{Conclusion}

The full framework can be summarized as
\begin{equation}
  \text{CR3BP physics}
  +
  \text{Bayesian estimation}
  +
  \text{adaptive feedback}.
\end{equation}
The CR3BP grounds the problem in nonlinear orbital mechanics. The EKF provides
the estimation engine. The navigation-risk score contributes an interpretable
feedback signal that connects covariance, geometry, missed data, and lighting
to adaptive filtering behavior.

Future work can strengthen the framework by deriving $\Phi_k$ explicitly,
defining $\GammaGeom(t)$ through the Fisher information matrix, connecting
$\Rnav(t)$ to entropy or information gain, and analyzing Kalman-gain stability.

\end{document}
